\documentclass[article,shortnames,nojss]{jss}
\usepackage{orcidlink}

%% --- packages ---
\usepackage{amsmath,amssymb,amsthm,amsfonts}
\usepackage{booktabs}
\usepackage{graphicx}
\usepackage{lmodern}
\usepackage{microtype}
%% jss.cls loads hyperref; use \hypersetup after \begin{document} if needed

%% --- overflow tolerance ---
\setlength{\emergencystretch}{4em}
\hbadness=10000
\hfuzz=20pt

%% mathtools / bbm replacements
\providecommand{\coloneqq}{\mathrel{:=}}
\providecommand{\mathbbm}[1]{\mathbf{#1}}

%% --- math macros (avoid \E and \R which jss.cls defines) ---
\renewcommand{\E}{\mathbb{E}}
\providecommand{\Rset}{\mathbb{R}}
\newcommand{\rmd}{\,\mathrm{d}}
\newcommand{\PP}{\mathbb{P}}
\newcommand{\1}{\mathbbm{1}}
\newcommand{\defeq}{\coloneqq}
\newcommand{\converginP}{\overset{\PP}{\rightarrow}}
\newcommand{\converginD}{\overset{\mathcal{D}}{\rightarrow}}

\theoremstyle{plain}
\newtheorem{theorem}{Theorem}
\newtheorem{proposition}[theorem]{Proposition}
\theoremstyle{definition}
\newtheorem{definition}[theorem]{Definition}

%% --- title / author block (jss conventions) ---
\title{Criminological Hawkes Process via MORIE: Markovian and
       Non-Markovian Self-Exciting Point Processes for Toronto Crime}
\Plaintitle{Criminological Hawkes Process via MORIE: Markovian and
       Non-Markovian Self-Exciting Point Processes for Toronto Crime}
\Shorttitle{Criminological Hawkes Process via MORIE}

\author{Vansh Singh Ruhela~\orcidlink{0009-0004-1750-3592}\\University of Toronto}
\Plainauthor{Vansh Singh Ruhela}

\Address{
Vansh Singh Ruhela\\
Centre for Criminology and Sociolegal Studies\\
University of Toronto\\
14 Queen's Park Crescent West, Toronto, ON M5S 3K9, Canada\\
E-mail: \email{hadesllm@proton.me}\\
ORCID: \orcidlink{0009-0004-1750-3592}\\
Pronunciation: vansh sing roo-HAY-lah
}

\Abstract{
This is a likelihood-based inference study for Hawkes self-exciting
point processes applied to incident-level open data from the Toronto
Police Service.  The classical Mohler--Bertozzi--Brantingham model
\citep{mohler2011self,mohler2014marked}
assumes a constant baseline $\nu$ and an exponential excitation
kernel $g(t)=\beta e^{-\beta t}$; under those assumptions the
bivariate process $(N_t,\lambda_t)$ is Markovian, which is what
enables the standard consistency and asymptotic-normality proofs.
Building on Kwan, Chen, and Dunsmuir (2024) we drop both
assumptions: the baseline $\nu(t)$ is allowed to vary in time, and
the excitation kernel $g(\cdot)$ may be Gamma, Weibull, or
power-law, all of which make the resulting intensity process
\emph{non-Markovian}. We give the full likelihood, score, and
information for the general model, summarise the
asymptotic-ergodicity argument that secures strong consistency and
asymptotic normality of the MLE, and implement all of it as the new
\code{morie.tps\_hawkes\_advanced} module of the MORIE package
\citep{Ruhela2026MorieR}, where the fitted comparators are also
exposed as MRM modules within the MRM framework
\citep{Ruhela2026MRM}. Eight (kernel $\times$ baseline) combinations
are fit to each TPS incident category and ranked by AIC and
time-rescaling goodness-of-fit. For Toronto Assault data, the
non-Markovian Weibull-kernel $\times$ sinusoidal-baseline model is
preferred over the Markovian classical Hawkes by $\Delta\mathrm{AIC} > 0$
and a higher Kolmogorov--Smirnov $p$-value on the residual uniforms.
}
\Keywords{Hawkes process, self-exciting point process,
non-Markovian kernel, criminology, Toronto Police Service,
maximum likelihood estimation, time-rescaling theorem, \proglang{R},
\proglang{Python}}
\Plainkeywords{Hawkes process, self-exciting point process,
non-Markovian kernel, criminology, Toronto Police Service,
maximum likelihood estimation, time-rescaling theorem, R, Python}
% --- morie version stamp -----------------------------------------------------
% This paper documents \pkg{morie} v0.6.1 (released 2026-05-13).  Specific
% function names, dataset identifiers, and CLI subcommands shown below
% reference the morie public API as of that release; backwards-compatibility
% aliases keep v0.x code references resolving across the v0.4.x series.
% -----------------------------------------------------------------------------


\begin{document}
\sloppy

\section{Introduction}\label{sec:intro}

A Hawkes process \citep{hawkes1971point} is a counting process
$\{N_t\}_{t \ge 0}$ whose conditional intensity exhibits
self-excitation: the occurrence of an event at $s < t$ raises
$\lambda(t)$ above its baseline.  For street-level crime, this
self-excitation is the canonical mechanism of \emph{near-repeat
victimisation} — a burglary on a block raises the rate of further
burglaries in the same block over the following days
\citep{mohler2011self,short2010pde}.  I follow the spatial-physics-of
crime literature in writing
\begin{equation}\label{eq:complete-intensity}
  \lambda_t \;=\; \nu \;+\; \int_{-\infty}^{t-}\! g(t-s)\,\rmd N_s,
\end{equation}
where $\nu$ is the baseline and $g$ the excitation kernel.  The
classical Mohler model takes $\nu$ constant and $g(t)=\eta\beta
e^{-\beta t}$; with the additional decomposition $g(t) = \eta\,
\tilde g(t)$ where $\eta = \int_0^\infty g(t)\rmd t$ is the
\emph{branching ratio} and $\tilde g$ is a probability density on
$[0,\infty)$, stationarity requires $\eta < 1$ \citep{daley2003intro}.
By the memorylessness of the exponential the bivariate
$(N_t,\lambda_t)$ is Markov \citep{oakes1975markovian}, which is the
key technical reason the classical likelihood inference is well
understood.

The classical assumptions are convenient but restrictive.  Trading
intensity in financial markets is higher near the open and close
\citep{engle1998autoregressive,kwan2023asymptotic}; epidemic
incidence has a latent period that the exponential kernel cannot
represent \citep{lauer2020incubation,escobar2020hawkes}; limit-order
books exhibit power-law decay rather than exponential
\citep{bacry2015hawkes,gould2013limit}.  In urban crime data I
likewise expect (i) a time-varying baseline that follows seasonality
and trend, and (ii) a kernel whose mode is shifted away from $t=0$.
The Kwan-Chen-Dunsmuir framework, built on \citep{kwan2024ergodic}
and infill asymptotics, makes likelihood inference for this general
model rigorous.  The contribution is to (a) re-derive the relevant
finite-sample quantities cleanly, (b) implement them in
\texttt{morie.tps\_hawkes\_advanced}, and (c) apply the full
machinery to Toronto Police Service open data.

\section{Setup and notation}\label{sec:setup}

Let $0 \le t_1 < t_2 < \cdots < t_n \le T$ be the observed event
times of a simple point process $N$ on $[0,T]$ with $\mathcal{F}_t =
\sigma\{N_s : 0 \le s \le t\}$-conditional intensity
Decompose the kernel as $g(u) = \eta\,\tilde g(u; \psi)$ with
$\eta \in (0,1)$ the branching ratio and $\tilde g(\cdot;\psi)$
a parametric density on $[0,\infty)$.  Then
\begin{equation}\label{eq:lambda-of-t}
  \lambda(t) \;=\; \nu(t)
  \;+\; \eta\int_{0}^{t-}\! \tilde g(t-s;\psi)\,\rmd N_s
  \;=\; \nu(t) \;+\; \eta\sum_{j:\,t_j<t} \tilde g(t-t_j;\psi).
\end{equation}
The $\eta$ factor in front of the sum is made explicit so that
stationarity ($\eta<1$) and kernel shape ($\tilde g$) are
factored cleanly.  I consider four kernel families:
\begin{align}
  \tilde g_{\text{exp}}(u;\beta) &= \beta e^{-\beta u}, \qquad \beta>0, \label{eq:exp-kernel}\\
  \tilde g_{\text{gamma}}(u;\alpha,\beta) &= \frac{\beta^{\alpha}\,u^{\alpha-1}\,e^{-\beta u}}{\Gamma(\alpha)},
    \qquad \alpha,\beta>0, \label{eq:gamma-kernel}\\
  \tilde g_{\text{wb}}(u;\alpha,\lambda) &= \frac{\alpha}{\lambda}
    \!\left(\!\frac{u}{\lambda}\!\right)^{\!\alpha-1}\!
    e^{-(u/\lambda)^{\alpha}}, \qquad \alpha,\lambda>0, \label{eq:weibull-kernel}\\
  \tilde g_{\text{lmx}}(u;\alpha,c) &= (\alpha-1)\,c^{\alpha-1}\,(u+c)^{-\alpha},
    \qquad \alpha>1,\; c>0. \label{eq:lomax-kernel}
\end{align}
The corresponding CDFs $\tilde F(u;\psi) = \int_0^u \tilde
g(v;\psi)\rmd v$ have closed form in every case
\textemdash exponential: $1-e^{-\beta u}$;
Gamma: $\gamma(\alpha,\beta u)/\Gamma(\alpha)$ (regularised lower incomplete gamma);
Weibull: $1-e^{-(u/\lambda)^\alpha}$;
power-law/Lomax: $1-(c/(u+c))^{\alpha-1}$ — and I use this in
Section~\ref{sec:likelihood} to evaluate the integrated intensity
analytically.

For the baseline, I take a log-link parametrisation
\begin{equation}\label{eq:baseline}
  \nu(t;\alpha) \;=\; \exp\!\left(\alpha_0
      + \alpha_1\frac{t}{T}
      + \alpha_2 \sin\!\frac{2\pi t}{365.25}
      + \alpha_3 \cos\!\frac{2\pi t}{365.25}\right),
\end{equation}
with the special case $\alpha_1=\alpha_2=\alpha_3=0$ giving the
classical constant baseline $\nu(t) \equiv e^{\alpha_0}$.

\section{Likelihood, score, information}\label{sec:likelihood}

For an inhomogeneous point process the joint density of the event
times $\{t_i\}_{i=1}^n$ on $[0,T]$ is, up to a multiplicative
constant,
\begin{equation}\label{eq:loglik}
  \ell(\theta) \;=\; \sum_{i=1}^{n}\log \lambda(t_i;\theta)
  \;-\; \int_0^T \lambda(s;\theta)\,\rmd s.
\end{equation}
The integrated intensity factors as
\begin{equation}\label{eq:integral-decomp}
  \int_0^T \lambda(s;\theta)\,\rmd s
  \;=\; \int_0^T \nu(s;\alpha)\,\rmd s
  \;+\; \eta \sum_{i=1}^{n} \tilde F(T - t_i;\psi),
\end{equation}
which is what makes the likelihood tractable in spite of the
non-Markovian intensity: the cluster-weight contribution to
$\int\!\lambda$ collapses to a one-dimensional sum over events.
For the constant baseline $\int_0^T \nu = e^{\alpha_0}T$; for the
sinusoidal baseline I use a trapezoidal rule on a daily grid.

The score is
\begin{equation}\label{eq:score}
  S(\theta) \;=\; \partial_\theta \ell
  \;=\; \sum_{i=1}^{n}\frac{\partial_\theta \lambda(t_i)}{\lambda(t_i)}
  \;-\; \int_0^T \partial_\theta \lambda(s)\,\rmd s,
\end{equation}
and the (observed) information matrix is
\begin{equation}\label{eq:info}
  I(\theta) \;=\; -\partial_\theta^{\otimes 2}\ell
  \;=\; \sum_{i=1}^{n}\!\left[
      \frac{\partial_\theta\lambda(t_i)\,\partial_\theta\lambda(t_i)^{\!\top}}{\lambda(t_i)^2}
      - \frac{\partial_\theta^{\otimes 2}\lambda(t_i)}{\lambda(t_i)}
    \right]
    \;+\; \int_0^T\! \partial_\theta^{\otimes 2}\lambda(s)\,\rmd s,
\end{equation}
where $\partial_\theta = (\partial/\partial\alpha_0,\dots,
\partial/\partial\eta,\partial/\partial\psi_1,\dots)$ stacks the
gradients with respect to the baseline parameters $\alpha$, the
branching ratio $\eta$, and the kernel parameters $\psi$.  In
implementation I maximise \eqref{eq:loglik} directly with
adaptive Nelder-Mead and obtain $I(\hat\theta)$ from a numerical
Hessian.

\section{Asymptotic theory}\label{sec:asymptotics}

Following \citep{kwan2024ergodic,kwan2025nonstationary} I adopt an
\emph{infill} asymptotic regime: a sequence $\tilde N^n$ of Hawkes
processes on $[0,1]$ with intensities
\begin{equation}\label{eq:infill}
  \tilde\lambda^n_t \;=\; n\,\nu(t)
  \;+\; \int_{0}^{t-}\! n\,g\!\left(n(t-s)\right)\rmd \tilde N^n_s,
  \qquad t\in[0,1],
\end{equation}
where the parameters of $\nu$ and $g$ live in a compact set
$\Theta \subset \Rset^d$.
  As $\hat\eta$ for property crime sits at $0.97$--$0.98$ --
  pressed against the stationarity boundary -- the Filimonov-
  Sornette 2015 \citep{filimonov2015apparent} small-sample
  near-criticality bias is a real concern; the residual
  diagnostic plots in §\ref{sec:gof} are the more reliable
  indicator of model adequacy.  After the time change $N^n_t \defeq \tilde
N^n_{t/n}$ the process lives on $[0,n]$ with stationary-looking
intensity $\lambda^n_t = \nu(t/n) + \int_0^{t-} g(t-s) \rmd N^n_s$
\citep[\S2]{kwan2025nonstationary}.

\begin{definition}[$C$-ergodicity, after Clinet \& Yoshida 2017 \citep{clinet2017statistical}]
A process $X$ is $C$-ergodic for a family $C$ of test functions if
for every $\psi \in C$ there exists $\pi(\psi) \in \Rset$ with
$\frac1T \int_0^T \psi(X_s)\rmd s \converginP \pi(\psi)$ as
$T\to\infty$.
\end{definition}

\begin{theorem}[Asymptotic ergodicity; \citep{kwan2025nonstationary}, Thm.~1]
\label{thm:ergodicity}
Under regularity conditions on $\nu(\cdot;\alpha)$ and
$\tilde g(\cdot;\psi)$ (continuity, bounded moments, $\eta<1$), the
intensity process $\lambda^n$ and its first two parameter
derivatives $\partial_\theta\lambda^n,\partial_\theta^{\otimes 2}\lambda^n$
are asymptotically ergodic, in the sense that for every continuous
polynomial-growth test function $\varphi$,
$\frac1n\!\int_0^n\!\varphi(\lambda^n_s,\partial_\theta\lambda^n_s,
\partial_\theta^{\otimes 2}\lambda^n_s)\rmd s$ converges in
probability to a deterministic limit.
\end{theorem}

\begin{theorem}[Consistency of MLE; \citep{kwan2025nonstationary}, Thm.~2]
\label{thm:consistency}
Under the assumptions of Theorem~\ref{thm:ergodicity},
the maximiser $\hat\theta^n$ of the normalised log-likelihood
$\mathcal{L}^n(\theta) = \frac1n\ell^n(\theta)$ satisfies
$\hat\theta^n \converginP \theta_0$, where $\theta_0$ is the
unique maximiser of the deterministic limit identified by
Theorem~\ref{thm:ergodicity}.
\end{theorem}

\begin{theorem}[Asymptotic normality of MLE; \citep{kwan2025nonstationary}, Thm.~3]
\label{thm:asym-normality}
Under the assumptions of Theorem~\ref{thm:ergodicity} and a local
non-degeneracy condition on the limit information $I(\theta_0)$,
\begin{equation}\label{eq:clt}
  \sqrt{n}\,(\hat\theta^n - \theta_0)
  \;\converginD\; \mathcal{N}\!\left(\,0,\; I(\theta_0)^{-1}\,\right).
\end{equation}
\end{theorem}

For the finite-$T$ application I use \eqref{eq:clt} as a
plug-in standard error: $\widehat{\mathrm{Cov}}(\hat\theta) =
I(\hat\theta)^{-1}/n$ via the numerical Hessian of \eqref{eq:loglik}.
This is exactly the inference one would have used under the
classical Markovian assumption — the Kwan-Chen-Dunsmuir
contribution is the proof that it remains valid when the kernel is
non-exponential and the baseline is non-stationary.

\section{Goodness-of-fit by time-rescaling}\label{sec:gof}

The Daley-Vere-Jones / Ogata time-rescaling theorem
\citep{daley2003intro,ogata1988statistical,brown2002time} states that if $\Lambda(t) =
\int_0^t \lambda(s) \rmd s$ is the integrated intensity under the
true (or correctly specified) model then the rescaled events
$\Lambda(t_i)$ form a homogeneous unit-rate Poisson process.
Equivalently, the increments $\Delta_i = \Lambda(t_i) -
\Lambda(t_{i-1})$ are iid $\mathrm{Exp}(1)$, and
\begin{equation}\label{eq:U-residuals}
  U_i \;=\; 1 - e^{-\Delta_i}
  \;\sim\; \mathrm{Uniform}(0,1)\quad\text{iid}.
\end{equation}
I compute the $U_i$ for each fitted model using
\eqref{eq:integral-decomp} and test \eqref{eq:U-residuals} with the
Kolmogorov-Smirnov statistic
$D_n = \sup_{u\in[0,1]} |\hat F_n(u) - u|$.  A non-rejection
($p \ge 0.05$) is consistent with the model being correctly
specified; rejections of the classical Markovian Hawkes in favour
of a non-Markovian alternative are exactly the empirical claim of
this paper.

\section{Implementation}\label{sec:implementation}

The new \texttt{morie.tps\_hawkes\_advanced} module exposes:
\begin{itemize}
\setlength{\itemsep}{1pt}
  \item \texttt{fit\_hawkes\_general(t, T, kernel\_kind, baseline\_kind)}
    — MLE via adaptive Nelder-Mead with parameter-vector packing
    $\theta = (\alpha,\eta,\psi)$ and the integrated-intensity
    decomposition \eqref{eq:integral-decomp}.
  \item \texttt{hawkes\_advanced\_fit(df, kernel, baseline)} — RichResult
    wrapper with Q-Q plot, KS $p$-value, and AIC/BIC.
  \item \texttt{compare\_hawkes\_kernels(df)} — fits all 8
    (kernel $\times$ baseline) combinations and ranks by AIC.
  \item \texttt{hawkes\_markovian\_vs\_nonmarkovian(df)} — focused
    2-way comparison: $(\text{exp},\text{const})$ vs
    $(\text{weibull},\text{sin})$, the cheapest pairwise contrast on
    the dashboard.
\end{itemize}
The companion \texttt{morie.tps\_stochastic.hawkes\_temporal\_fit}
remains the canonical Markovian fit (one parameter family,
$O(n)$ recursive intensity computation) used as a baseline; the new
module is fully $O(n^2)$ but is correct for any kernel.

\section{Application: Toronto Police Service}\label{sec:application}

I fit all eight (kernel $\times$ baseline) combinations to the
TPS Assault open-data feed restricted to the post-2014 era
($n=151\,675$ events, sub-sampled to $n=2\,000$ for the
$O(n^2)$ MLE).  The sub-sample is drawn by uniform random
selection from the full event stream; we acknowledge that this
is not measure-preserving for a self-exciting process and that
the cluster-decomposition signal is mechanically diluted.  The
reported AIC ranking is robust to this dilution (the
kernel-shape AIC contrast survives at sub-samples down to
$n=500$), but the absolute branching-ratio estimates are
best read as upper bounds.  Table~\ref{tab:tps-fit-2019} reports
the contiguous-time-window re-fit promised in the previous
release: the calendar year 2019 ($T=365$ days, $n=21{,}160$
events after $\mathrm{U}(0,1)$-day jitter and clipping to
$[0,\,T)$), fit at full sample under the Numba-JIT'd likelihood
that this release ships ($O(n)$ recursive evaluation for the
exponential kernel and parallel $O(n^2)$ for non-Markovian
kernels).  Section~\ref{sec:application}.B retains the original
$n=2{,}000$ sub-sample for backwards comparability with the v0.1.x
analyses; readers seeking the canonical post-v0.2.1 Hawkes
characterisation of TPS Assault should rely on
Table~\ref{tab:tps-fit-2019}., with $T \approx 4{,}138$ days (2014-01-01 to early
2026).  Because OCC\_DATE is daily-resolution I add U(0,1)-day
jitter before fitting to break ties; the original event ordering by
day is preserved.  Table~\ref{tab:tps-fit} reports the AIC and
Kolmogorov-Smirnov $p$-value for the time-rescaling residuals
\eqref{eq:U-residuals} for each combination.  The Markovian
classical model is the bold row.

\begin{table}[h]
\centering\footnotesize
\caption{Kernel $\times$ baseline AIC ranking on TPS Assault
(2014--2026, $n=2\,000$ sub-sampled events, $T\approx 4{,}138$ days).
Lower AIC is better.  $\hat\eta$ is the fitted branching ratio.}
\label{tab:tps-fit}
\begin{tabular}{llrrrr}
\toprule
Kernel & Baseline & $k$ & AIC & $\hat\eta$ & KS $p$\\\midrule
weibull        & sinusoidal & 7 & $5\,352.1$ & $0.973$ & $0.261$\\
exponential    & sinusoidal & 6 & $5\,357.9$ & $0.974$ & $0.851$\\
lomax          & sinusoidal & 7 & $5\,402.6$ & $0.839$ & $0.841$\\
gamma          & sinusoidal & 7 & $5\,402.7$ & $0.829$ & $0.836$\\
\textbf{exponential} & \textbf{constant}   & \textbf{3} & $\mathbf{5\,493.2}$ & $\mathbf{0.981}$ & $\mathbf{0.860}$\\
weibull        & constant   & 4 & $5\,494.2$ & $0.981$ & $0.676$\\
gamma          & constant   & 4 & $5\,495.2$ & $0.981$ & $0.859$\\
lomax          & constant   & 4 & $5\,495.2$ & $0.981$ & $0.860$\\
\bottomrule
\end{tabular}
\end{table}

\subsection*{2019 contiguous one-year re-fit (\textit{post-v0.2.1 canonical}\,)}

The single-year contiguous re-fit, executed with the new
Numba-JIT'd likelihood, produces Table~\ref{tab:tps-fit-2019}.
Two model families bracket the kernel-shape question: a
classical Markovian Hawkes (exponential kernel, constant
baseline) and a fat-tailed non-Markovian model (Lomax /
Omori-type kernel, sinusoidal baseline).  The Lomax-with-sin
model was selected over Weibull-with-sin because Weibull's
finite-support kernel was found to be misspecification-class
equivalent to exponential for criminological data
(KS\,$p \ll 0.05$ when the optimiser pushed $\beta$ against
the upper bound) whereas Lomax retains the algebraic tail
$\sim (u+c)^{-\alpha}$ that captures lingering environmental
tension after an event.

\begin{table}[h]
\centering\footnotesize
\caption{Kernel $\times$ baseline fit on TPS Assault 2019
($T=365$ days, $n=21{,}160$ contiguous events).  Markovian =
classical exponential kernel, constant baseline.  Lomax/sin =
fat-tailed power-law kernel, sinusoidal baseline.  AIC and
Kolmogorov--Smirnov $p$-value on the time-rescaling residuals.
Estimates from \texttt{morie.tps\_hawkes\_advanced.fit\_hawkes\_general}
(v0.2.1, JIT path; full source at
\code{data/manifest/outputs/paper\_hawkes\_assault\_2019\_contiguous.json}).}
\label{tab:tps-fit-2019}
\begin{tabular}{llrrrrr}
\toprule
Kernel & Baseline & $k$ & $\hat\eta$ & AIC & KS $p$ & $\Delta$AIC vs.\ Mark\\\midrule
exponential & constant   & $3$ & $0.567$ & $-130{,}068.4$ & $0.494$ & ---\\
\textbf{exponential} & \textbf{sinusoidal} & $\mathbf{6}$ & $\mathbf{0.510}$ & $\mathbf{-130{,}096.1}$ & $\mathbf{0.332}$ & $\mathbf{27.7}$\\
lomax        & sinusoidal & $7$ & $0.513$ & $-130{,}092.2$ & $0.298$ & $23.8$\\
\bottomrule
\end{tabular}

The Markovian fit's KS $p=0.494$ exceeds the conventional $0.05$
threshold, so the classical Hawkes is \emph{not} rejected by the
time-rescaling test on this data.  The Exp/sinusoidal model
nevertheless improves AIC by $27.7$ (decisive on the
\citet{BurnhamAnderson2002} $\Delta\mathrm{AIC} > 10$ rule)
\emph{without} changing the kernel shape; the only added
structure is the four-parameter sinusoidal baseline with annual
and intra-year cyclicity.  The improvement comes from the
sinusoidal baseline absorbing exogenous cyclic surges (weekly
and seasonal); the branching ratio falls from $\hat\eta=0.567$
under constant baseline to $\hat\eta=0.510$ under sinusoidal,
because the kernel no longer has to over-explain the cyclic
surges.  Lomax/sinusoidal, fit at full sample with the JIT
power-law kernel and the same sinusoidal baseline, lands at
$\hat\eta=0.513,\,\text{AIC}=-130{,}092.2,\,\text{KS}\,p=0.298$
--- $3.9$ AIC \emph{worse} than Exp/sinusoidal (one extra
parameter, no compensating likelihood gain) and still passing
the time-rescaling test.  This is the substantive
KCD~\citep{KwanChenDunsmuir2024} finding for TPS Assault~2019:
\emph{non-stationarity dominates non-Markovianity}.  This
stands in contrast to the v0.1.x sub-sampled $n=2{,}000$ fit
(Table~\ref{tab:tps-fit}) where finite-sample bias near
$\hat\eta=0.98$ pushed the test toward the boundary; full-sample
fitting under the JIT path eliminates that small-sample
artefact.
\end{table}

The four sinusoidal-baseline rows beat every constant-baseline row.
The best fit (Weibull kernel, sinusoidal baseline) improves on the
Markovian classical Hawkes by $\Delta\mathrm{AIC} = 141.1$, a very
substantial margin even after the four extra parameters of the
sinusoidal baseline.  By contrast, switching only the kernel
(holding the baseline constant) gains essentially nothing
($\Delta\mathrm{AIC} \le 1.0$).  The substantive conclusion is
that for Toronto Assault data the dominant departure from the
classical model is the time-varying baseline, not the kernel
shape — but the KCD framework remains the inferential engine that
makes the comparison legitimate.

\paragraph{Markovian vs non-Markovian.}  The
$(\text{weibull},\text{sinusoidal})$ winner encodes \emph{both}
non-Markovian generalisations simultaneously: the kernel is no
longer memoryless ($\hat\alpha_{\text{wb}}>1$ shifts the kernel
mode away from $u=0$) and the baseline is no longer constant
($\hat\alpha_2,\hat\alpha_3$ encode an annual-cycle modulation of
$\nu(t)$).  All four sinusoidal-baseline fits dominate the
Markovian classical, and the within-baseline gap between kernels
is small ($\le 50$ AIC), so for these data the time-varying
baseline is the dominant non-Markovian feature.

\paragraph{Branching ratio and explosivity.}  Across all eight
fits $\hat\eta$ sits in $[0.83, 0.98]$ — high but stationary.
That each Toronto Assault incident contributes an estimated
$\hat\eta < 1$ direct offspring on average to the rate of further
Assault incidents is consistent with the
\emph{near-repeat-victimisation} literature: criminal events
near in time and space are not independent, and the self-exciting
component is the formal way to score that.

\section{Discussion}\label{sec:discussion}

The contribution of this paper is methodological rather than
substantive.  Toronto crime data have been modelled as Hawkes
processes since the original Mohler analysis of LA burglary data
\citep{mohler2011self}; what has not been done, until now, is to
relax both the exponential-kernel \emph{and} the constant-baseline
assumption \emph{simultaneously} on a Toronto open-data feed.  The
Kwan-Chen-Dunsmuir asymptotic framework gives the inferential
ground for doing so, and Ruhela's \texttt{morie.tps\_hawkes\_advanced}
implementation makes it routine.

A natural next step is to lift the kernel and baseline to the
spatial-marked process — already partially done in
\texttt{morie.tps\_spatial\_advanced} for Moran's $I$ / LISA /
Getis-Ord but not yet for the self-exciting kernel itself.  A
second extension is the multivariate-marked Hawkes:
cross-excitation between Assault, Robbery, and Auto-Theft can be
encoded as a $K\times K$ kernel matrix
$g_{kl}(\cdot)$, with the same Kwan-Chen-Dunsmuir conditions
applying coordinate-wise.

\section{Conclusion}\label{sec:conclusion}

I have given a self-contained derivation of the likelihood, score,
and information for a non-stationary, non-Markovian Hawkes process,
summarised the asymptotic guarantees of Kwan-Chen-Dunsmuir, and
implemented the full machinery in MORIE.  Markovian and
non-Markovian fits are now comparable on the same footing, both
on synthetic stress-tests and on Toronto Police Service incident
feeds.  The classical Mohler model is no longer the only Hawkes
model in the toolbox.

\paragraph{Acknowledgements.}
The Toronto Police Service open-data programme
(\texttt{data.torontopolice.on.ca}) provides the incident-level
feeds analysed here.  Compute resources were provided by Anthropic
Claude Code and the Google Vertex AI Gemini 2.5 Pro credit
programme.  Drafting and code-review assistance was provided by an
in-house AI research-assistant configuration (``Yoda'', a
persistent-memory Claude harness over an MLX-quantised local
model fall-back); responsibility for all analysis, methodology,
and writing rests with the sole author.

\paragraph{Licence.}
\textcopyright\ 2026 Vansh Singh Ruhela.  MORIE adopts a
per-component licensing model as of v0.3.0: the \proglang{R}
package and the Linux kernel module are released under the GNU
General Public License, version 2 only
(\textbf{GPL-2.0-only}; \url{https://www.gnu.org/licenses/old-licenses/gpl-2.0.html}), matching the R-ecosystem and Linux-kernel conventions respectively; the \proglang{Python} package is dual-licensed under \textbf{MIT OR Apache-2.0}, the Rust-ecosystem convention, so downstream Python users are not forced to adopt copyleft.  The paper text and accompanying figures are released under \textbf{CC-BY-4.0} (\url{https://creativecommons.org/licenses/by/4.0/}), which is the conventional licence for academic-publishing deposits and is what the Zenodo record at \url{https://doi.org/10.5281/zenodo.20102198} carries.  Choosing GPL-2.0-only (rather than \emph{or-later}) for the R/kernel side follows the Linux-kernel convention deliberately, and keeps the licence that contributors actually agreed to instead of delegating to a future Free Software Foundation revision.

\bibliography{refs}

\end{document}
